\documentclass[11pt,a4paper]{article}
% main (standalone preamble for editing)
% vim: nowrap: spell:
% ══════════════════════════════════════════════════════════════════════════════
% Speed Maths --- shared preamble
% Included by every sheet and answer file via:  % main (standalone preamble for editing)
% vim: nowrap: spell:
% ══════════════════════════════════════════════════════════════════════════════
% Speed Maths --- shared preamble
% Included by every sheet and answer file via:  % main (standalone preamble for editing)
% vim: nowrap: spell:
% ══════════════════════════════════════════════════════════════════════════════
% Speed Maths --- shared preamble
% Included by every sheet and answer file via:  \input{../../shared/preamble}
% Editing THIS file changes the look of every sheet/answer across every pillar.
% ══════════════════════════════════════════════════════════════════════════════

\usepackage[a4paper, top=25mm, bottom=25mm, left=22mm, right=22mm]{geometry}
\usepackage{amsmath, amssymb}
\usepackage{enumitem}
\usepackage{booktabs}
\usepackage{array}
\usepackage{parskip}
\usepackage{microtype}
\usepackage{fancyhdr}
\usepackage{titlesec}
\usepackage{mdframed}
\usepackage{xcolor}
\usepackage[hidelinks]{hyperref}
\usepackage{graphicx}
\usepackage{tikz}
\usepackage{pgfplots}
\pgfplotsset{compat=1.18}

% ── Colours ───────────────────────────────────────────────────────────────────
\definecolor{darkgray}{gray}{0.15}
\definecolor{medgray}{gray}{0.45}
\definecolor{lightgray}{gray}{0.93}
\definecolor{boxgray}{gray}{0.88}

% ── Section formatting ──────────────────────────────────────────────────────────
\titleformat{\section}[block]
  {\normalfont\large\bfseries}{}
  {0pt}{}[\vspace{2pt}\hrule\vspace{6pt}]
\titlespacing*{\section}{0pt}{14pt}{4pt}

% ── List formatting ─────────────────────────────────────────────────────────────
\setlist[enumerate,1]{label=\textbf{\arabic*.}, leftmargin=2em, itemsep=10pt}

% ── Branding constants (edit here, not per-file) ────────────────────────────────
\newcommand{\SpeedExamLine}{\textbf{Speed Maths:} Competition \& Exam Prep}
\newcommand{\SpeedCredit}{Series by \href{https://www.linkedin.com/in/cerealdev/}{David Akinyele-Aje}}

% ── Header / footer ──────────────────────────────────────────────────────────────
% Usage (once, after \input{../../shared/preamble} and before \begin{document}):
%   \SpeedHeader{Algebra}{3}
% First arg  = pillar name shown as "Speed <Pillar> --- Daily Drill"
% Second arg = worksheet number
\pagestyle{fancy}
\newcommand{\SpeedHeader}[2]{%
  \fancyhf{}%
  \renewcommand{\headrulewidth}{0.4pt}%
  \setlength{\headheight}{14pt}%
  \fancyhead[L]{\small\textbf{Speed #1 --- Daily Drill}}%
  \fancyhead[R]{\small Worksheet \##2}%
  \fancyfoot[L]{\small \SpeedExamLine}%
  \fancyfoot[C]{\small\thepage}%
  \fancyfoot[R]{\small \SpeedCredit}%
  \renewcommand{\footrulewidth}{0.4pt}%
}

% ── Title block (used at the top of every sheet AND answer file) ───────────────
% Usage: \SpeedTitleBlock{Daily Algebra Drill \#3}{Jane Doe}
% %% AGENTS: When generating a new sheet, ask the human contributor for the name they want credited in argument 2.
% For answer files, pass the "--- Answers \& Investigations" suffix in the arg.
\newcommand{\SpeedTitleBlock}[2]{%
  \begin{center}
    {\LARGE\bfseries #1}\\[4pt]
    {\normalsize\itshape Sheet by #2}\\[6pt]
    \hrule\vspace{4pt}
    \begin{tabular}{llcc}
      \toprule
      \textbf{Section} & \textbf{Topic} & \textbf{Questions} & \textbf{Target time}\\
      \midrule
      A & Rapid Recognition          & 10 & 2 min 30 sec \\
      B & Manipulation Drills        & 10 & 8 min        \\
      C & Substitution \& Structure  & 8  & 10 min       \\
      D & Challenge                  & 5  & 15 min       \\
      \bottomrule
    \end{tabular}
    \vspace{4pt}
    \hrule
  \end{center}
}

% ── Sheet metadata: topic + the day's new tools ────────────────────────────────
% Usage (immediately after \SpeedTitleBlock, sheets only):
%   \SpeedMeta{Expanding, factorising, and the identity toolkit}
%             {difference of squares, sum/product form, completing the square}
%
% Arg 1 = the sheet's topic in one line. The website lists this instead of
%         "Sheet 03", so write the thing a reader would actually search for.
% Arg 2 = comma-separated, at most five: the tools this sheet introduces that
%         earlier sheets in the pillar did not. These become the website's tags.
%         Leave out anything the reader already met — the tags are meant to say
%         what is new today, not to summarise the sheet.
%
% Both are parsed straight out of this file by tools/build_website.py, so the
% sheet stays the single source of truth and the site cannot drift from it.
%
% This renders NOTHING. It is a declaration for the build script, not a piece of
% the sheet's layout: adding it to a sheet must not change that sheet's PDF, so
% the 25 existing sheets can be annotated without a single one of them being
% reissued. If the toolkit should also appear on the page, that is a separate
% decision about the sheet design — make it deliberately, here, not by leaving
% this macro printing something nobody asked it to.
\newcommand{\SpeedMeta}[2]{}

% ── Answer / Method / Investigate-further boxes (answer files only) ────────────
\newmdenv[
  backgroundcolor=lightgray,
  linecolor=medgray,
  linewidth=0.5pt,
  leftmargin=0pt, rightmargin=0pt,
  innerleftmargin=8pt, innerrightmargin=8pt,
  innertopmargin=5pt, innerbottommargin=5pt,
  skipabove=4pt, skipbelow=4pt
]{ansbox}

\newmdenv[
  backgroundcolor=white,
  linecolor=darkgray,
  linewidth=0.8pt,
  leftline=true, rightline=false, topline=false, bottomline=false,
  leftmargin=0pt, rightmargin=0pt,
  innerleftmargin=8pt, innerrightmargin=6pt,
  innertopmargin=4pt, innerbottommargin=4pt,
  skipabove=4pt, skipbelow=6pt
]{invbox}

\newcommand{\ans}[1]{%
  \begin{ansbox}
    \textbf{Answer:} #1
  \end{ansbox}}

\newcommand{\method}[1]{%
  {\small\textbf{Method:} #1}\par}

\newcommand{\inv}[1]{%
  \begin{invbox}
    \textbf{\small Investigate further:} {\small #1}
  \end{invbox}}

% ── Misc helpers ────────────────────────────────────────────────────────────────
\newcommand{\qn}[1]{\textbf{#1.}\;}

% Mistake-linking: attach a Top 5 Patterns / Common Traps bullet to the question
% label(s) in this sheet where it actually shows up, e.g. \seealso{B6, D2}
\newcommand{\seealso}[1]{\hfill\textit{\footnotesize(see #1)}}

% ── Graph options macro ─────────────────────────────────────────────────────────
% Usage: \GraphOptions{tikz code for A}{tikz code for B}{tikz code for C}{tikz code for D}
\newcommand{\GraphOptions}[4]{%
  \begin{center}
    \begin{tabular}{cc}
      \textbf{A)} & \textbf{B)} \\
      #1 & #2 \\[10pt]
      \textbf{C)} & \textbf{D)} \\
      #3 & #4
    \end{tabular}
  \end{center}
}

% ── Closing motivational rule + cross-promo footer (sheets only) ────────────────
% Usage: \SpeedClosing{"Quote text."}
\newcommand{\SpeedClosing}[1]{%
  \vspace{20pt}
  \begin{center}
  \rule{0.6\textwidth}{0.4pt}\\[4pt]
  {\small\itshape #1}\\[4pt]
  \rule{0.6\textwidth}{0.4pt}
  \end{center}
  \begin{center}
  {\small Found a mistake or want to contribute a sheet?}\\
  {\small Visit the open-source project at \href{https://github.com/The-CerealDev/speed-maths}{github.com/The-CerealDev/speed-maths}}
  \end{center}
}

% Editing THIS file changes the look of every sheet/answer across every pillar.
% ══════════════════════════════════════════════════════════════════════════════

\usepackage[a4paper, top=25mm, bottom=25mm, left=22mm, right=22mm]{geometry}
\usepackage{amsmath, amssymb}
\usepackage{enumitem}
\usepackage{booktabs}
\usepackage{array}
\usepackage{parskip}
\usepackage{microtype}
\usepackage{fancyhdr}
\usepackage{titlesec}
\usepackage{mdframed}
\usepackage{xcolor}
\usepackage[hidelinks]{hyperref}
\usepackage{graphicx}
\usepackage{tikz}
\usepackage{pgfplots}
\pgfplotsset{compat=1.18}

% ── Colours ───────────────────────────────────────────────────────────────────
\definecolor{darkgray}{gray}{0.15}
\definecolor{medgray}{gray}{0.45}
\definecolor{lightgray}{gray}{0.93}
\definecolor{boxgray}{gray}{0.88}

% ── Section formatting ──────────────────────────────────────────────────────────
\titleformat{\section}[block]
  {\normalfont\large\bfseries}{}
  {0pt}{}[\vspace{2pt}\hrule\vspace{6pt}]
\titlespacing*{\section}{0pt}{14pt}{4pt}

% ── List formatting ─────────────────────────────────────────────────────────────
\setlist[enumerate,1]{label=\textbf{\arabic*.}, leftmargin=2em, itemsep=10pt}

% ── Branding constants (edit here, not per-file) ────────────────────────────────
\newcommand{\SpeedExamLine}{\textbf{Speed Maths:} Competition \& Exam Prep}
\newcommand{\SpeedCredit}{Series by \href{https://www.linkedin.com/in/cerealdev/}{David Akinyele-Aje}}

% ── Header / footer ──────────────────────────────────────────────────────────────
% Usage (once, after % main (standalone preamble for editing)
% vim: nowrap: spell:
% ══════════════════════════════════════════════════════════════════════════════
% Speed Maths --- shared preamble
% Included by every sheet and answer file via:  \input{../../shared/preamble}
% Editing THIS file changes the look of every sheet/answer across every pillar.
% ══════════════════════════════════════════════════════════════════════════════

\usepackage[a4paper, top=25mm, bottom=25mm, left=22mm, right=22mm]{geometry}
\usepackage{amsmath, amssymb}
\usepackage{enumitem}
\usepackage{booktabs}
\usepackage{array}
\usepackage{parskip}
\usepackage{microtype}
\usepackage{fancyhdr}
\usepackage{titlesec}
\usepackage{mdframed}
\usepackage{xcolor}
\usepackage[hidelinks]{hyperref}
\usepackage{graphicx}
\usepackage{tikz}
\usepackage{pgfplots}
\pgfplotsset{compat=1.18}

% ── Colours ───────────────────────────────────────────────────────────────────
\definecolor{darkgray}{gray}{0.15}
\definecolor{medgray}{gray}{0.45}
\definecolor{lightgray}{gray}{0.93}
\definecolor{boxgray}{gray}{0.88}

% ── Section formatting ──────────────────────────────────────────────────────────
\titleformat{\section}[block]
  {\normalfont\large\bfseries}{}
  {0pt}{}[\vspace{2pt}\hrule\vspace{6pt}]
\titlespacing*{\section}{0pt}{14pt}{4pt}

% ── List formatting ─────────────────────────────────────────────────────────────
\setlist[enumerate,1]{label=\textbf{\arabic*.}, leftmargin=2em, itemsep=10pt}

% ── Branding constants (edit here, not per-file) ────────────────────────────────
\newcommand{\SpeedExamLine}{\textbf{Speed Maths:} Competition \& Exam Prep}
\newcommand{\SpeedCredit}{Series by \href{https://www.linkedin.com/in/cerealdev/}{David Akinyele-Aje}}

% ── Header / footer ──────────────────────────────────────────────────────────────
% Usage (once, after \input{../../shared/preamble} and before \begin{document}):
%   \SpeedHeader{Algebra}{3}
% First arg  = pillar name shown as "Speed <Pillar> --- Daily Drill"
% Second arg = worksheet number
\pagestyle{fancy}
\newcommand{\SpeedHeader}[2]{%
  \fancyhf{}%
  \renewcommand{\headrulewidth}{0.4pt}%
  \setlength{\headheight}{14pt}%
  \fancyhead[L]{\small\textbf{Speed #1 --- Daily Drill}}%
  \fancyhead[R]{\small Worksheet \##2}%
  \fancyfoot[L]{\small \SpeedExamLine}%
  \fancyfoot[C]{\small\thepage}%
  \fancyfoot[R]{\small \SpeedCredit}%
  \renewcommand{\footrulewidth}{0.4pt}%
}

% ── Title block (used at the top of every sheet AND answer file) ───────────────
% Usage: \SpeedTitleBlock{Daily Algebra Drill \#3}{Jane Doe}
% %% AGENTS: When generating a new sheet, ask the human contributor for the name they want credited in argument 2.
% For answer files, pass the "--- Answers \& Investigations" suffix in the arg.
\newcommand{\SpeedTitleBlock}[2]{%
  \begin{center}
    {\LARGE\bfseries #1}\\[4pt]
    {\normalsize\itshape Sheet by #2}\\[6pt]
    \hrule\vspace{4pt}
    \begin{tabular}{llcc}
      \toprule
      \textbf{Section} & \textbf{Topic} & \textbf{Questions} & \textbf{Target time}\\
      \midrule
      A & Rapid Recognition          & 10 & 2 min 30 sec \\
      B & Manipulation Drills        & 10 & 8 min        \\
      C & Substitution \& Structure  & 8  & 10 min       \\
      D & Challenge                  & 5  & 15 min       \\
      \bottomrule
    \end{tabular}
    \vspace{4pt}
    \hrule
  \end{center}
}

% ── Sheet metadata: topic + the day's new tools ────────────────────────────────
% Usage (immediately after \SpeedTitleBlock, sheets only):
%   \SpeedMeta{Expanding, factorising, and the identity toolkit}
%             {difference of squares, sum/product form, completing the square}
%
% Arg 1 = the sheet's topic in one line. The website lists this instead of
%         "Sheet 03", so write the thing a reader would actually search for.
% Arg 2 = comma-separated, at most five: the tools this sheet introduces that
%         earlier sheets in the pillar did not. These become the website's tags.
%         Leave out anything the reader already met — the tags are meant to say
%         what is new today, not to summarise the sheet.
%
% Both are parsed straight out of this file by tools/build_website.py, so the
% sheet stays the single source of truth and the site cannot drift from it.
%
% This renders NOTHING. It is a declaration for the build script, not a piece of
% the sheet's layout: adding it to a sheet must not change that sheet's PDF, so
% the 25 existing sheets can be annotated without a single one of them being
% reissued. If the toolkit should also appear on the page, that is a separate
% decision about the sheet design — make it deliberately, here, not by leaving
% this macro printing something nobody asked it to.
\newcommand{\SpeedMeta}[2]{}

% ── Answer / Method / Investigate-further boxes (answer files only) ────────────
\newmdenv[
  backgroundcolor=lightgray,
  linecolor=medgray,
  linewidth=0.5pt,
  leftmargin=0pt, rightmargin=0pt,
  innerleftmargin=8pt, innerrightmargin=8pt,
  innertopmargin=5pt, innerbottommargin=5pt,
  skipabove=4pt, skipbelow=4pt
]{ansbox}

\newmdenv[
  backgroundcolor=white,
  linecolor=darkgray,
  linewidth=0.8pt,
  leftline=true, rightline=false, topline=false, bottomline=false,
  leftmargin=0pt, rightmargin=0pt,
  innerleftmargin=8pt, innerrightmargin=6pt,
  innertopmargin=4pt, innerbottommargin=4pt,
  skipabove=4pt, skipbelow=6pt
]{invbox}

\newcommand{\ans}[1]{%
  \begin{ansbox}
    \textbf{Answer:} #1
  \end{ansbox}}

\newcommand{\method}[1]{%
  {\small\textbf{Method:} #1}\par}

\newcommand{\inv}[1]{%
  \begin{invbox}
    \textbf{\small Investigate further:} {\small #1}
  \end{invbox}}

% ── Misc helpers ────────────────────────────────────────────────────────────────
\newcommand{\qn}[1]{\textbf{#1.}\;}

% Mistake-linking: attach a Top 5 Patterns / Common Traps bullet to the question
% label(s) in this sheet where it actually shows up, e.g. \seealso{B6, D2}
\newcommand{\seealso}[1]{\hfill\textit{\footnotesize(see #1)}}

% ── Graph options macro ─────────────────────────────────────────────────────────
% Usage: \GraphOptions{tikz code for A}{tikz code for B}{tikz code for C}{tikz code for D}
\newcommand{\GraphOptions}[4]{%
  \begin{center}
    \begin{tabular}{cc}
      \textbf{A)} & \textbf{B)} \\
      #1 & #2 \\[10pt]
      \textbf{C)} & \textbf{D)} \\
      #3 & #4
    \end{tabular}
  \end{center}
}

% ── Closing motivational rule + cross-promo footer (sheets only) ────────────────
% Usage: \SpeedClosing{"Quote text."}
\newcommand{\SpeedClosing}[1]{%
  \vspace{20pt}
  \begin{center}
  \rule{0.6\textwidth}{0.4pt}\\[4pt]
  {\small\itshape #1}\\[4pt]
  \rule{0.6\textwidth}{0.4pt}
  \end{center}
  \begin{center}
  {\small Found a mistake or want to contribute a sheet?}\\
  {\small Visit the open-source project at \href{https://github.com/The-CerealDev/speed-maths}{github.com/The-CerealDev/speed-maths}}
  \end{center}
}
 and before \begin{document}):
%   \SpeedHeader{Algebra}{3}
% First arg  = pillar name shown as "Speed <Pillar> --- Daily Drill"
% Second arg = worksheet number
\pagestyle{fancy}
\newcommand{\SpeedHeader}[2]{%
  \fancyhf{}%
  \renewcommand{\headrulewidth}{0.4pt}%
  \setlength{\headheight}{14pt}%
  \fancyhead[L]{\small\textbf{Speed #1 --- Daily Drill}}%
  \fancyhead[R]{\small Worksheet \##2}%
  \fancyfoot[L]{\small \SpeedExamLine}%
  \fancyfoot[C]{\small\thepage}%
  \fancyfoot[R]{\small \SpeedCredit}%
  \renewcommand{\footrulewidth}{0.4pt}%
}

% ── Title block (used at the top of every sheet AND answer file) ───────────────
% Usage: \SpeedTitleBlock{Daily Algebra Drill \#3}{Jane Doe}
% %% AGENTS: When generating a new sheet, ask the human contributor for the name they want credited in argument 2.
% For answer files, pass the "--- Answers \& Investigations" suffix in the arg.
\newcommand{\SpeedTitleBlock}[2]{%
  \begin{center}
    {\LARGE\bfseries #1}\\[4pt]
    {\normalsize\itshape Sheet by #2}\\[6pt]
    \hrule\vspace{4pt}
    \begin{tabular}{llcc}
      \toprule
      \textbf{Section} & \textbf{Topic} & \textbf{Questions} & \textbf{Target time}\\
      \midrule
      A & Rapid Recognition          & 10 & 2 min 30 sec \\
      B & Manipulation Drills        & 10 & 8 min        \\
      C & Substitution \& Structure  & 8  & 10 min       \\
      D & Challenge                  & 5  & 15 min       \\
      \bottomrule
    \end{tabular}
    \vspace{4pt}
    \hrule
  \end{center}
}

% ── Sheet metadata: topic + the day's new tools ────────────────────────────────
% Usage (immediately after \SpeedTitleBlock, sheets only):
%   \SpeedMeta{Expanding, factorising, and the identity toolkit}
%             {difference of squares, sum/product form, completing the square}
%
% Arg 1 = the sheet's topic in one line. The website lists this instead of
%         "Sheet 03", so write the thing a reader would actually search for.
% Arg 2 = comma-separated, at most five: the tools this sheet introduces that
%         earlier sheets in the pillar did not. These become the website's tags.
%         Leave out anything the reader already met — the tags are meant to say
%         what is new today, not to summarise the sheet.
%
% Both are parsed straight out of this file by tools/build_website.py, so the
% sheet stays the single source of truth and the site cannot drift from it.
%
% This renders NOTHING. It is a declaration for the build script, not a piece of
% the sheet's layout: adding it to a sheet must not change that sheet's PDF, so
% the 25 existing sheets can be annotated without a single one of them being
% reissued. If the toolkit should also appear on the page, that is a separate
% decision about the sheet design — make it deliberately, here, not by leaving
% this macro printing something nobody asked it to.
\newcommand{\SpeedMeta}[2]{}

% ── Answer / Method / Investigate-further boxes (answer files only) ────────────
\newmdenv[
  backgroundcolor=lightgray,
  linecolor=medgray,
  linewidth=0.5pt,
  leftmargin=0pt, rightmargin=0pt,
  innerleftmargin=8pt, innerrightmargin=8pt,
  innertopmargin=5pt, innerbottommargin=5pt,
  skipabove=4pt, skipbelow=4pt
]{ansbox}

\newmdenv[
  backgroundcolor=white,
  linecolor=darkgray,
  linewidth=0.8pt,
  leftline=true, rightline=false, topline=false, bottomline=false,
  leftmargin=0pt, rightmargin=0pt,
  innerleftmargin=8pt, innerrightmargin=6pt,
  innertopmargin=4pt, innerbottommargin=4pt,
  skipabove=4pt, skipbelow=6pt
]{invbox}

\newcommand{\ans}[1]{%
  \begin{ansbox}
    \textbf{Answer:} #1
  \end{ansbox}}

\newcommand{\method}[1]{%
  {\small\textbf{Method:} #1}\par}

\newcommand{\inv}[1]{%
  \begin{invbox}
    \textbf{\small Investigate further:} {\small #1}
  \end{invbox}}

% ── Misc helpers ────────────────────────────────────────────────────────────────
\newcommand{\qn}[1]{\textbf{#1.}\;}

% Mistake-linking: attach a Top 5 Patterns / Common Traps bullet to the question
% label(s) in this sheet where it actually shows up, e.g. \seealso{B6, D2}
\newcommand{\seealso}[1]{\hfill\textit{\footnotesize(see #1)}}

% ── Graph options macro ─────────────────────────────────────────────────────────
% Usage: \GraphOptions{tikz code for A}{tikz code for B}{tikz code for C}{tikz code for D}
\newcommand{\GraphOptions}[4]{%
  \begin{center}
    \begin{tabular}{cc}
      \textbf{A)} & \textbf{B)} \\
      #1 & #2 \\[10pt]
      \textbf{C)} & \textbf{D)} \\
      #3 & #4
    \end{tabular}
  \end{center}
}

% ── Closing motivational rule + cross-promo footer (sheets only) ────────────────
% Usage: \SpeedClosing{"Quote text."}
\newcommand{\SpeedClosing}[1]{%
  \vspace{20pt}
  \begin{center}
  \rule{0.6\textwidth}{0.4pt}\\[4pt]
  {\small\itshape #1}\\[4pt]
  \rule{0.6\textwidth}{0.4pt}
  \end{center}
  \begin{center}
  {\small Found a mistake or want to contribute a sheet?}\\
  {\small Visit the open-source project at \href{https://github.com/The-CerealDev/speed-maths}{github.com/The-CerealDev/speed-maths}}
  \end{center}
}

% Editing THIS file changes the look of every sheet/answer across every pillar.
% ══════════════════════════════════════════════════════════════════════════════

\usepackage[a4paper, top=25mm, bottom=25mm, left=22mm, right=22mm]{geometry}
\usepackage{amsmath, amssymb}
\usepackage{enumitem}
\usepackage{booktabs}
\usepackage{array}
\usepackage{parskip}
\usepackage{microtype}
\usepackage{fancyhdr}
\usepackage{titlesec}
\usepackage{mdframed}
\usepackage{xcolor}
\usepackage[hidelinks]{hyperref}
\usepackage{graphicx}
\usepackage{tikz}
\usepackage{pgfplots}
\pgfplotsset{compat=1.18}

% ── Colours ───────────────────────────────────────────────────────────────────
\definecolor{darkgray}{gray}{0.15}
\definecolor{medgray}{gray}{0.45}
\definecolor{lightgray}{gray}{0.93}
\definecolor{boxgray}{gray}{0.88}

% ── Section formatting ──────────────────────────────────────────────────────────
\titleformat{\section}[block]
  {\normalfont\large\bfseries}{}
  {0pt}{}[\vspace{2pt}\hrule\vspace{6pt}]
\titlespacing*{\section}{0pt}{14pt}{4pt}

% ── List formatting ─────────────────────────────────────────────────────────────
\setlist[enumerate,1]{label=\textbf{\arabic*.}, leftmargin=2em, itemsep=10pt}

% ── Branding constants (edit here, not per-file) ────────────────────────────────
\newcommand{\SpeedExamLine}{\textbf{Speed Maths:} Competition \& Exam Prep}
\newcommand{\SpeedCredit}{Series by \href{https://www.linkedin.com/in/cerealdev/}{David Akinyele-Aje}}

% ── Header / footer ──────────────────────────────────────────────────────────────
% Usage (once, after % main (standalone preamble for editing)
% vim: nowrap: spell:
% ══════════════════════════════════════════════════════════════════════════════
% Speed Maths --- shared preamble
% Included by every sheet and answer file via:  % main (standalone preamble for editing)
% vim: nowrap: spell:
% ══════════════════════════════════════════════════════════════════════════════
% Speed Maths --- shared preamble
% Included by every sheet and answer file via:  \input{../../shared/preamble}
% Editing THIS file changes the look of every sheet/answer across every pillar.
% ══════════════════════════════════════════════════════════════════════════════

\usepackage[a4paper, top=25mm, bottom=25mm, left=22mm, right=22mm]{geometry}
\usepackage{amsmath, amssymb}
\usepackage{enumitem}
\usepackage{booktabs}
\usepackage{array}
\usepackage{parskip}
\usepackage{microtype}
\usepackage{fancyhdr}
\usepackage{titlesec}
\usepackage{mdframed}
\usepackage{xcolor}
\usepackage[hidelinks]{hyperref}
\usepackage{graphicx}
\usepackage{tikz}
\usepackage{pgfplots}
\pgfplotsset{compat=1.18}

% ── Colours ───────────────────────────────────────────────────────────────────
\definecolor{darkgray}{gray}{0.15}
\definecolor{medgray}{gray}{0.45}
\definecolor{lightgray}{gray}{0.93}
\definecolor{boxgray}{gray}{0.88}

% ── Section formatting ──────────────────────────────────────────────────────────
\titleformat{\section}[block]
  {\normalfont\large\bfseries}{}
  {0pt}{}[\vspace{2pt}\hrule\vspace{6pt}]
\titlespacing*{\section}{0pt}{14pt}{4pt}

% ── List formatting ─────────────────────────────────────────────────────────────
\setlist[enumerate,1]{label=\textbf{\arabic*.}, leftmargin=2em, itemsep=10pt}

% ── Branding constants (edit here, not per-file) ────────────────────────────────
\newcommand{\SpeedExamLine}{\textbf{Speed Maths:} Competition \& Exam Prep}
\newcommand{\SpeedCredit}{Series by \href{https://www.linkedin.com/in/cerealdev/}{David Akinyele-Aje}}

% ── Header / footer ──────────────────────────────────────────────────────────────
% Usage (once, after \input{../../shared/preamble} and before \begin{document}):
%   \SpeedHeader{Algebra}{3}
% First arg  = pillar name shown as "Speed <Pillar> --- Daily Drill"
% Second arg = worksheet number
\pagestyle{fancy}
\newcommand{\SpeedHeader}[2]{%
  \fancyhf{}%
  \renewcommand{\headrulewidth}{0.4pt}%
  \setlength{\headheight}{14pt}%
  \fancyhead[L]{\small\textbf{Speed #1 --- Daily Drill}}%
  \fancyhead[R]{\small Worksheet \##2}%
  \fancyfoot[L]{\small \SpeedExamLine}%
  \fancyfoot[C]{\small\thepage}%
  \fancyfoot[R]{\small \SpeedCredit}%
  \renewcommand{\footrulewidth}{0.4pt}%
}

% ── Title block (used at the top of every sheet AND answer file) ───────────────
% Usage: \SpeedTitleBlock{Daily Algebra Drill \#3}{Jane Doe}
% %% AGENTS: When generating a new sheet, ask the human contributor for the name they want credited in argument 2.
% For answer files, pass the "--- Answers \& Investigations" suffix in the arg.
\newcommand{\SpeedTitleBlock}[2]{%
  \begin{center}
    {\LARGE\bfseries #1}\\[4pt]
    {\normalsize\itshape Sheet by #2}\\[6pt]
    \hrule\vspace{4pt}
    \begin{tabular}{llcc}
      \toprule
      \textbf{Section} & \textbf{Topic} & \textbf{Questions} & \textbf{Target time}\\
      \midrule
      A & Rapid Recognition          & 10 & 2 min 30 sec \\
      B & Manipulation Drills        & 10 & 8 min        \\
      C & Substitution \& Structure  & 8  & 10 min       \\
      D & Challenge                  & 5  & 15 min       \\
      \bottomrule
    \end{tabular}
    \vspace{4pt}
    \hrule
  \end{center}
}

% ── Sheet metadata: topic + the day's new tools ────────────────────────────────
% Usage (immediately after \SpeedTitleBlock, sheets only):
%   \SpeedMeta{Expanding, factorising, and the identity toolkit}
%             {difference of squares, sum/product form, completing the square}
%
% Arg 1 = the sheet's topic in one line. The website lists this instead of
%         "Sheet 03", so write the thing a reader would actually search for.
% Arg 2 = comma-separated, at most five: the tools this sheet introduces that
%         earlier sheets in the pillar did not. These become the website's tags.
%         Leave out anything the reader already met — the tags are meant to say
%         what is new today, not to summarise the sheet.
%
% Both are parsed straight out of this file by tools/build_website.py, so the
% sheet stays the single source of truth and the site cannot drift from it.
%
% This renders NOTHING. It is a declaration for the build script, not a piece of
% the sheet's layout: adding it to a sheet must not change that sheet's PDF, so
% the 25 existing sheets can be annotated without a single one of them being
% reissued. If the toolkit should also appear on the page, that is a separate
% decision about the sheet design — make it deliberately, here, not by leaving
% this macro printing something nobody asked it to.
\newcommand{\SpeedMeta}[2]{}

% ── Answer / Method / Investigate-further boxes (answer files only) ────────────
\newmdenv[
  backgroundcolor=lightgray,
  linecolor=medgray,
  linewidth=0.5pt,
  leftmargin=0pt, rightmargin=0pt,
  innerleftmargin=8pt, innerrightmargin=8pt,
  innertopmargin=5pt, innerbottommargin=5pt,
  skipabove=4pt, skipbelow=4pt
]{ansbox}

\newmdenv[
  backgroundcolor=white,
  linecolor=darkgray,
  linewidth=0.8pt,
  leftline=true, rightline=false, topline=false, bottomline=false,
  leftmargin=0pt, rightmargin=0pt,
  innerleftmargin=8pt, innerrightmargin=6pt,
  innertopmargin=4pt, innerbottommargin=4pt,
  skipabove=4pt, skipbelow=6pt
]{invbox}

\newcommand{\ans}[1]{%
  \begin{ansbox}
    \textbf{Answer:} #1
  \end{ansbox}}

\newcommand{\method}[1]{%
  {\small\textbf{Method:} #1}\par}

\newcommand{\inv}[1]{%
  \begin{invbox}
    \textbf{\small Investigate further:} {\small #1}
  \end{invbox}}

% ── Misc helpers ────────────────────────────────────────────────────────────────
\newcommand{\qn}[1]{\textbf{#1.}\;}

% Mistake-linking: attach a Top 5 Patterns / Common Traps bullet to the question
% label(s) in this sheet where it actually shows up, e.g. \seealso{B6, D2}
\newcommand{\seealso}[1]{\hfill\textit{\footnotesize(see #1)}}

% ── Graph options macro ─────────────────────────────────────────────────────────
% Usage: \GraphOptions{tikz code for A}{tikz code for B}{tikz code for C}{tikz code for D}
\newcommand{\GraphOptions}[4]{%
  \begin{center}
    \begin{tabular}{cc}
      \textbf{A)} & \textbf{B)} \\
      #1 & #2 \\[10pt]
      \textbf{C)} & \textbf{D)} \\
      #3 & #4
    \end{tabular}
  \end{center}
}

% ── Closing motivational rule + cross-promo footer (sheets only) ────────────────
% Usage: \SpeedClosing{"Quote text."}
\newcommand{\SpeedClosing}[1]{%
  \vspace{20pt}
  \begin{center}
  \rule{0.6\textwidth}{0.4pt}\\[4pt]
  {\small\itshape #1}\\[4pt]
  \rule{0.6\textwidth}{0.4pt}
  \end{center}
  \begin{center}
  {\small Found a mistake or want to contribute a sheet?}\\
  {\small Visit the open-source project at \href{https://github.com/The-CerealDev/speed-maths}{github.com/The-CerealDev/speed-maths}}
  \end{center}
}

% Editing THIS file changes the look of every sheet/answer across every pillar.
% ══════════════════════════════════════════════════════════════════════════════

\usepackage[a4paper, top=25mm, bottom=25mm, left=22mm, right=22mm]{geometry}
\usepackage{amsmath, amssymb}
\usepackage{enumitem}
\usepackage{booktabs}
\usepackage{array}
\usepackage{parskip}
\usepackage{microtype}
\usepackage{fancyhdr}
\usepackage{titlesec}
\usepackage{mdframed}
\usepackage{xcolor}
\usepackage[hidelinks]{hyperref}
\usepackage{graphicx}
\usepackage{tikz}
\usepackage{pgfplots}
\pgfplotsset{compat=1.18}

% ── Colours ───────────────────────────────────────────────────────────────────
\definecolor{darkgray}{gray}{0.15}
\definecolor{medgray}{gray}{0.45}
\definecolor{lightgray}{gray}{0.93}
\definecolor{boxgray}{gray}{0.88}

% ── Section formatting ──────────────────────────────────────────────────────────
\titleformat{\section}[block]
  {\normalfont\large\bfseries}{}
  {0pt}{}[\vspace{2pt}\hrule\vspace{6pt}]
\titlespacing*{\section}{0pt}{14pt}{4pt}

% ── List formatting ─────────────────────────────────────────────────────────────
\setlist[enumerate,1]{label=\textbf{\arabic*.}, leftmargin=2em, itemsep=10pt}

% ── Branding constants (edit here, not per-file) ────────────────────────────────
\newcommand{\SpeedExamLine}{\textbf{Speed Maths:} Competition \& Exam Prep}
\newcommand{\SpeedCredit}{Series by \href{https://www.linkedin.com/in/cerealdev/}{David Akinyele-Aje}}

% ── Header / footer ──────────────────────────────────────────────────────────────
% Usage (once, after % main (standalone preamble for editing)
% vim: nowrap: spell:
% ══════════════════════════════════════════════════════════════════════════════
% Speed Maths --- shared preamble
% Included by every sheet and answer file via:  \input{../../shared/preamble}
% Editing THIS file changes the look of every sheet/answer across every pillar.
% ══════════════════════════════════════════════════════════════════════════════

\usepackage[a4paper, top=25mm, bottom=25mm, left=22mm, right=22mm]{geometry}
\usepackage{amsmath, amssymb}
\usepackage{enumitem}
\usepackage{booktabs}
\usepackage{array}
\usepackage{parskip}
\usepackage{microtype}
\usepackage{fancyhdr}
\usepackage{titlesec}
\usepackage{mdframed}
\usepackage{xcolor}
\usepackage[hidelinks]{hyperref}
\usepackage{graphicx}
\usepackage{tikz}
\usepackage{pgfplots}
\pgfplotsset{compat=1.18}

% ── Colours ───────────────────────────────────────────────────────────────────
\definecolor{darkgray}{gray}{0.15}
\definecolor{medgray}{gray}{0.45}
\definecolor{lightgray}{gray}{0.93}
\definecolor{boxgray}{gray}{0.88}

% ── Section formatting ──────────────────────────────────────────────────────────
\titleformat{\section}[block]
  {\normalfont\large\bfseries}{}
  {0pt}{}[\vspace{2pt}\hrule\vspace{6pt}]
\titlespacing*{\section}{0pt}{14pt}{4pt}

% ── List formatting ─────────────────────────────────────────────────────────────
\setlist[enumerate,1]{label=\textbf{\arabic*.}, leftmargin=2em, itemsep=10pt}

% ── Branding constants (edit here, not per-file) ────────────────────────────────
\newcommand{\SpeedExamLine}{\textbf{Speed Maths:} Competition \& Exam Prep}
\newcommand{\SpeedCredit}{Series by \href{https://www.linkedin.com/in/cerealdev/}{David Akinyele-Aje}}

% ── Header / footer ──────────────────────────────────────────────────────────────
% Usage (once, after \input{../../shared/preamble} and before \begin{document}):
%   \SpeedHeader{Algebra}{3}
% First arg  = pillar name shown as "Speed <Pillar> --- Daily Drill"
% Second arg = worksheet number
\pagestyle{fancy}
\newcommand{\SpeedHeader}[2]{%
  \fancyhf{}%
  \renewcommand{\headrulewidth}{0.4pt}%
  \setlength{\headheight}{14pt}%
  \fancyhead[L]{\small\textbf{Speed #1 --- Daily Drill}}%
  \fancyhead[R]{\small Worksheet \##2}%
  \fancyfoot[L]{\small \SpeedExamLine}%
  \fancyfoot[C]{\small\thepage}%
  \fancyfoot[R]{\small \SpeedCredit}%
  \renewcommand{\footrulewidth}{0.4pt}%
}

% ── Title block (used at the top of every sheet AND answer file) ───────────────
% Usage: \SpeedTitleBlock{Daily Algebra Drill \#3}{Jane Doe}
% %% AGENTS: When generating a new sheet, ask the human contributor for the name they want credited in argument 2.
% For answer files, pass the "--- Answers \& Investigations" suffix in the arg.
\newcommand{\SpeedTitleBlock}[2]{%
  \begin{center}
    {\LARGE\bfseries #1}\\[4pt]
    {\normalsize\itshape Sheet by #2}\\[6pt]
    \hrule\vspace{4pt}
    \begin{tabular}{llcc}
      \toprule
      \textbf{Section} & \textbf{Topic} & \textbf{Questions} & \textbf{Target time}\\
      \midrule
      A & Rapid Recognition          & 10 & 2 min 30 sec \\
      B & Manipulation Drills        & 10 & 8 min        \\
      C & Substitution \& Structure  & 8  & 10 min       \\
      D & Challenge                  & 5  & 15 min       \\
      \bottomrule
    \end{tabular}
    \vspace{4pt}
    \hrule
  \end{center}
}

% ── Sheet metadata: topic + the day's new tools ────────────────────────────────
% Usage (immediately after \SpeedTitleBlock, sheets only):
%   \SpeedMeta{Expanding, factorising, and the identity toolkit}
%             {difference of squares, sum/product form, completing the square}
%
% Arg 1 = the sheet's topic in one line. The website lists this instead of
%         "Sheet 03", so write the thing a reader would actually search for.
% Arg 2 = comma-separated, at most five: the tools this sheet introduces that
%         earlier sheets in the pillar did not. These become the website's tags.
%         Leave out anything the reader already met — the tags are meant to say
%         what is new today, not to summarise the sheet.
%
% Both are parsed straight out of this file by tools/build_website.py, so the
% sheet stays the single source of truth and the site cannot drift from it.
%
% This renders NOTHING. It is a declaration for the build script, not a piece of
% the sheet's layout: adding it to a sheet must not change that sheet's PDF, so
% the 25 existing sheets can be annotated without a single one of them being
% reissued. If the toolkit should also appear on the page, that is a separate
% decision about the sheet design — make it deliberately, here, not by leaving
% this macro printing something nobody asked it to.
\newcommand{\SpeedMeta}[2]{}

% ── Answer / Method / Investigate-further boxes (answer files only) ────────────
\newmdenv[
  backgroundcolor=lightgray,
  linecolor=medgray,
  linewidth=0.5pt,
  leftmargin=0pt, rightmargin=0pt,
  innerleftmargin=8pt, innerrightmargin=8pt,
  innertopmargin=5pt, innerbottommargin=5pt,
  skipabove=4pt, skipbelow=4pt
]{ansbox}

\newmdenv[
  backgroundcolor=white,
  linecolor=darkgray,
  linewidth=0.8pt,
  leftline=true, rightline=false, topline=false, bottomline=false,
  leftmargin=0pt, rightmargin=0pt,
  innerleftmargin=8pt, innerrightmargin=6pt,
  innertopmargin=4pt, innerbottommargin=4pt,
  skipabove=4pt, skipbelow=6pt
]{invbox}

\newcommand{\ans}[1]{%
  \begin{ansbox}
    \textbf{Answer:} #1
  \end{ansbox}}

\newcommand{\method}[1]{%
  {\small\textbf{Method:} #1}\par}

\newcommand{\inv}[1]{%
  \begin{invbox}
    \textbf{\small Investigate further:} {\small #1}
  \end{invbox}}

% ── Misc helpers ────────────────────────────────────────────────────────────────
\newcommand{\qn}[1]{\textbf{#1.}\;}

% Mistake-linking: attach a Top 5 Patterns / Common Traps bullet to the question
% label(s) in this sheet where it actually shows up, e.g. \seealso{B6, D2}
\newcommand{\seealso}[1]{\hfill\textit{\footnotesize(see #1)}}

% ── Graph options macro ─────────────────────────────────────────────────────────
% Usage: \GraphOptions{tikz code for A}{tikz code for B}{tikz code for C}{tikz code for D}
\newcommand{\GraphOptions}[4]{%
  \begin{center}
    \begin{tabular}{cc}
      \textbf{A)} & \textbf{B)} \\
      #1 & #2 \\[10pt]
      \textbf{C)} & \textbf{D)} \\
      #3 & #4
    \end{tabular}
  \end{center}
}

% ── Closing motivational rule + cross-promo footer (sheets only) ────────────────
% Usage: \SpeedClosing{"Quote text."}
\newcommand{\SpeedClosing}[1]{%
  \vspace{20pt}
  \begin{center}
  \rule{0.6\textwidth}{0.4pt}\\[4pt]
  {\small\itshape #1}\\[4pt]
  \rule{0.6\textwidth}{0.4pt}
  \end{center}
  \begin{center}
  {\small Found a mistake or want to contribute a sheet?}\\
  {\small Visit the open-source project at \href{https://github.com/The-CerealDev/speed-maths}{github.com/The-CerealDev/speed-maths}}
  \end{center}
}
 and before \begin{document}):
%   \SpeedHeader{Algebra}{3}
% First arg  = pillar name shown as "Speed <Pillar> --- Daily Drill"
% Second arg = worksheet number
\pagestyle{fancy}
\newcommand{\SpeedHeader}[2]{%
  \fancyhf{}%
  \renewcommand{\headrulewidth}{0.4pt}%
  \setlength{\headheight}{14pt}%
  \fancyhead[L]{\small\textbf{Speed #1 --- Daily Drill}}%
  \fancyhead[R]{\small Worksheet \##2}%
  \fancyfoot[L]{\small \SpeedExamLine}%
  \fancyfoot[C]{\small\thepage}%
  \fancyfoot[R]{\small \SpeedCredit}%
  \renewcommand{\footrulewidth}{0.4pt}%
}

% ── Title block (used at the top of every sheet AND answer file) ───────────────
% Usage: \SpeedTitleBlock{Daily Algebra Drill \#3}{Jane Doe}
% %% AGENTS: When generating a new sheet, ask the human contributor for the name they want credited in argument 2.
% For answer files, pass the "--- Answers \& Investigations" suffix in the arg.
\newcommand{\SpeedTitleBlock}[2]{%
  \begin{center}
    {\LARGE\bfseries #1}\\[4pt]
    {\normalsize\itshape Sheet by #2}\\[6pt]
    \hrule\vspace{4pt}
    \begin{tabular}{llcc}
      \toprule
      \textbf{Section} & \textbf{Topic} & \textbf{Questions} & \textbf{Target time}\\
      \midrule
      A & Rapid Recognition          & 10 & 2 min 30 sec \\
      B & Manipulation Drills        & 10 & 8 min        \\
      C & Substitution \& Structure  & 8  & 10 min       \\
      D & Challenge                  & 5  & 15 min       \\
      \bottomrule
    \end{tabular}
    \vspace{4pt}
    \hrule
  \end{center}
}

% ── Sheet metadata: topic + the day's new tools ────────────────────────────────
% Usage (immediately after \SpeedTitleBlock, sheets only):
%   \SpeedMeta{Expanding, factorising, and the identity toolkit}
%             {difference of squares, sum/product form, completing the square}
%
% Arg 1 = the sheet's topic in one line. The website lists this instead of
%         "Sheet 03", so write the thing a reader would actually search for.
% Arg 2 = comma-separated, at most five: the tools this sheet introduces that
%         earlier sheets in the pillar did not. These become the website's tags.
%         Leave out anything the reader already met — the tags are meant to say
%         what is new today, not to summarise the sheet.
%
% Both are parsed straight out of this file by tools/build_website.py, so the
% sheet stays the single source of truth and the site cannot drift from it.
%
% This renders NOTHING. It is a declaration for the build script, not a piece of
% the sheet's layout: adding it to a sheet must not change that sheet's PDF, so
% the 25 existing sheets can be annotated without a single one of them being
% reissued. If the toolkit should also appear on the page, that is a separate
% decision about the sheet design — make it deliberately, here, not by leaving
% this macro printing something nobody asked it to.
\newcommand{\SpeedMeta}[2]{}

% ── Answer / Method / Investigate-further boxes (answer files only) ────────────
\newmdenv[
  backgroundcolor=lightgray,
  linecolor=medgray,
  linewidth=0.5pt,
  leftmargin=0pt, rightmargin=0pt,
  innerleftmargin=8pt, innerrightmargin=8pt,
  innertopmargin=5pt, innerbottommargin=5pt,
  skipabove=4pt, skipbelow=4pt
]{ansbox}

\newmdenv[
  backgroundcolor=white,
  linecolor=darkgray,
  linewidth=0.8pt,
  leftline=true, rightline=false, topline=false, bottomline=false,
  leftmargin=0pt, rightmargin=0pt,
  innerleftmargin=8pt, innerrightmargin=6pt,
  innertopmargin=4pt, innerbottommargin=4pt,
  skipabove=4pt, skipbelow=6pt
]{invbox}

\newcommand{\ans}[1]{%
  \begin{ansbox}
    \textbf{Answer:} #1
  \end{ansbox}}

\newcommand{\method}[1]{%
  {\small\textbf{Method:} #1}\par}

\newcommand{\inv}[1]{%
  \begin{invbox}
    \textbf{\small Investigate further:} {\small #1}
  \end{invbox}}

% ── Misc helpers ────────────────────────────────────────────────────────────────
\newcommand{\qn}[1]{\textbf{#1.}\;}

% Mistake-linking: attach a Top 5 Patterns / Common Traps bullet to the question
% label(s) in this sheet where it actually shows up, e.g. \seealso{B6, D2}
\newcommand{\seealso}[1]{\hfill\textit{\footnotesize(see #1)}}

% ── Graph options macro ─────────────────────────────────────────────────────────
% Usage: \GraphOptions{tikz code for A}{tikz code for B}{tikz code for C}{tikz code for D}
\newcommand{\GraphOptions}[4]{%
  \begin{center}
    \begin{tabular}{cc}
      \textbf{A)} & \textbf{B)} \\
      #1 & #2 \\[10pt]
      \textbf{C)} & \textbf{D)} \\
      #3 & #4
    \end{tabular}
  \end{center}
}

% ── Closing motivational rule + cross-promo footer (sheets only) ────────────────
% Usage: \SpeedClosing{"Quote text."}
\newcommand{\SpeedClosing}[1]{%
  \vspace{20pt}
  \begin{center}
  \rule{0.6\textwidth}{0.4pt}\\[4pt]
  {\small\itshape #1}\\[4pt]
  \rule{0.6\textwidth}{0.4pt}
  \end{center}
  \begin{center}
  {\small Found a mistake or want to contribute a sheet?}\\
  {\small Visit the open-source project at \href{https://github.com/The-CerealDev/speed-maths}{github.com/The-CerealDev/speed-maths}}
  \end{center}
}
 and before \begin{document}):
%   \SpeedHeader{Algebra}{3}
% First arg  = pillar name shown as "Speed <Pillar> --- Daily Drill"
% Second arg = worksheet number
\pagestyle{fancy}
\newcommand{\SpeedHeader}[2]{%
  \fancyhf{}%
  \renewcommand{\headrulewidth}{0.4pt}%
  \setlength{\headheight}{14pt}%
  \fancyhead[L]{\small\textbf{Speed #1 --- Daily Drill}}%
  \fancyhead[R]{\small Worksheet \##2}%
  \fancyfoot[L]{\small \SpeedExamLine}%
  \fancyfoot[C]{\small\thepage}%
  \fancyfoot[R]{\small \SpeedCredit}%
  \renewcommand{\footrulewidth}{0.4pt}%
}

% ── Title block (used at the top of every sheet AND answer file) ───────────────
% Usage: \SpeedTitleBlock{Daily Algebra Drill \#3}{Jane Doe}
% %% AGENTS: When generating a new sheet, ask the human contributor for the name they want credited in argument 2.
% For answer files, pass the "--- Answers \& Investigations" suffix in the arg.
\newcommand{\SpeedTitleBlock}[2]{%
  \begin{center}
    {\LARGE\bfseries #1}\\[4pt]
    {\normalsize\itshape Sheet by #2}\\[6pt]
    \hrule\vspace{4pt}
    \begin{tabular}{llcc}
      \toprule
      \textbf{Section} & \textbf{Topic} & \textbf{Questions} & \textbf{Target time}\\
      \midrule
      A & Rapid Recognition          & 10 & 2 min 30 sec \\
      B & Manipulation Drills        & 10 & 8 min        \\
      C & Substitution \& Structure  & 8  & 10 min       \\
      D & Challenge                  & 5  & 15 min       \\
      \bottomrule
    \end{tabular}
    \vspace{4pt}
    \hrule
  \end{center}
}

% ── Sheet metadata: topic + the day's new tools ────────────────────────────────
% Usage (immediately after \SpeedTitleBlock, sheets only):
%   \SpeedMeta{Expanding, factorising, and the identity toolkit}
%             {difference of squares, sum/product form, completing the square}
%
% Arg 1 = the sheet's topic in one line. The website lists this instead of
%         "Sheet 03", so write the thing a reader would actually search for.
% Arg 2 = comma-separated, at most five: the tools this sheet introduces that
%         earlier sheets in the pillar did not. These become the website's tags.
%         Leave out anything the reader already met — the tags are meant to say
%         what is new today, not to summarise the sheet.
%
% Both are parsed straight out of this file by tools/build_website.py, so the
% sheet stays the single source of truth and the site cannot drift from it.
%
% This renders NOTHING. It is a declaration for the build script, not a piece of
% the sheet's layout: adding it to a sheet must not change that sheet's PDF, so
% the 25 existing sheets can be annotated without a single one of them being
% reissued. If the toolkit should also appear on the page, that is a separate
% decision about the sheet design — make it deliberately, here, not by leaving
% this macro printing something nobody asked it to.
\newcommand{\SpeedMeta}[2]{}

% ── Answer / Method / Investigate-further boxes (answer files only) ────────────
\newmdenv[
  backgroundcolor=lightgray,
  linecolor=medgray,
  linewidth=0.5pt,
  leftmargin=0pt, rightmargin=0pt,
  innerleftmargin=8pt, innerrightmargin=8pt,
  innertopmargin=5pt, innerbottommargin=5pt,
  skipabove=4pt, skipbelow=4pt
]{ansbox}

\newmdenv[
  backgroundcolor=white,
  linecolor=darkgray,
  linewidth=0.8pt,
  leftline=true, rightline=false, topline=false, bottomline=false,
  leftmargin=0pt, rightmargin=0pt,
  innerleftmargin=8pt, innerrightmargin=6pt,
  innertopmargin=4pt, innerbottommargin=4pt,
  skipabove=4pt, skipbelow=6pt
]{invbox}

\newcommand{\ans}[1]{%
  \begin{ansbox}
    \textbf{Answer:} #1
  \end{ansbox}}

\newcommand{\method}[1]{%
  {\small\textbf{Method:} #1}\par}

\newcommand{\inv}[1]{%
  \begin{invbox}
    \textbf{\small Investigate further:} {\small #1}
  \end{invbox}}

% ── Misc helpers ────────────────────────────────────────────────────────────────
\newcommand{\qn}[1]{\textbf{#1.}\;}

% Mistake-linking: attach a Top 5 Patterns / Common Traps bullet to the question
% label(s) in this sheet where it actually shows up, e.g. \seealso{B6, D2}
\newcommand{\seealso}[1]{\hfill\textit{\footnotesize(see #1)}}

% ── Graph options macro ─────────────────────────────────────────────────────────
% Usage: \GraphOptions{tikz code for A}{tikz code for B}{tikz code for C}{tikz code for D}
\newcommand{\GraphOptions}[4]{%
  \begin{center}
    \begin{tabular}{cc}
      \textbf{A)} & \textbf{B)} \\
      #1 & #2 \\[10pt]
      \textbf{C)} & \textbf{D)} \\
      #3 & #4
    \end{tabular}
  \end{center}
}

% ── Closing motivational rule + cross-promo footer (sheets only) ────────────────
% Usage: \SpeedClosing{"Quote text."}
\newcommand{\SpeedClosing}[1]{%
  \vspace{20pt}
  \begin{center}
  \rule{0.6\textwidth}{0.4pt}\\[4pt]
  {\small\itshape #1}\\[4pt]
  \rule{0.6\textwidth}{0.4pt}
  \end{center}
  \begin{center}
  {\small Found a mistake or want to contribute a sheet?}\\
  {\small Visit the open-source project at \href{https://github.com/The-CerealDev/speed-maths}{github.com/The-CerealDev/speed-maths}}
  \end{center}
}

\SpeedHeader{Sequences}{2}

\begin{document}

\SpeedTitleBlock{Daily Sequences Drill \#2}{David Akinyele-Aje}
\noindent\textit{New toolkit today: Binomial \& Trinomial Expansion $\cdot$ Pascal's Triangle as a Sequence $\cdot$ The Generalized Binomial Series.}

\vspace{2pt}
\noindent\textit{Attempt each question before checking answers. Time yourself. No calculators.}

\section*{Section A \quad Rapid Recognition \hfill{\normalfont\small($\leq$15\,s each)}}
% ══════════════════════════════════════════════════════════════════════════════

\noindent\textit{Pascal's triangle and binomial coefficients on sight.}

\begin{enumerate}

\item Write down row $5$ of Pascal's triangle (the coefficients of $(x+y)^5$).

\item Evaluate $\dbinom{6}{2}$.

\item Write down $\displaystyle\sum_{k=0}^{n}\binom nk$ in closed form.

\item Expand $(1+x)^3$ fully.

\item Find the coefficient of $x^2$ in $(1+x)^6$.

\item Row $4$ of Pascal's triangle is $1,4,6,4,1$. Is this sequence of five numbers arithmetic? Justify with one check.

\item Simplify $\dbinom n0+\dbinom n1$.

\item State the value of $\displaystyle\sum_{k=0}^{n}(-1)^k\binom nk$ for $n\geq1$.

\item Find the constant term in $(1+x)^5$.

\item True or false: row $n$ of Pascal's triangle always has exactly $n+1$ entries.

\end{enumerate}

% ══════════════════════════════════════════════════════════════════════════════
\section*{Section B \quad Manipulation Drills \hfill{\normalfont\small($\leq$50\,s each)}}
% ══════════════════════════════════════════════════════════════════════════════

\noindent\textit{Substitute into the binomial expansion; extract specific coefficients.}

\begin{enumerate}

\item Prove that $\dbinom n0+\dbinom n1+\cdots+\dbinom nn=2^n$ by substituting a value into the expansion of $(1+x)^n$.

\item Row $10$ of Pascal's triangle sums to $2^{10}$. If instead you sum only its first four entries, $\binom{10}{0}+\binom{10}{1}+\binom{10}{2}+\binom{10}{3}$, is the result arithmetic, geometric, or neither, as a rule? Compute the value.

\item Using $x=-1$ in $(1+x)^n$ ($n\geq1$), what is $\displaystyle\sum_{k=0}^n(-1)^k\binom nk$?
\begin{itemize}
\item[A)] $0$
\item[B)] $1$
\item[C)] $2^n$
\item[D)] $(-1)^n$
\end{itemize}

\item Find the coefficient of $x^3$ in $(2+x)^5$.
\begin{itemize}
\item[A)] $10$
\item[B)] $40$
\item[C)] $80$
\item[D)] $20$
\end{itemize}

\item Find the coefficient of $x^4$ in $(1-2x)^6$.
\begin{itemize}
\item[A)] $60$
\item[B)] $-240$
\item[C)] $240$
\item[D)] $-60$
\end{itemize}

\item Find the constant term in the expansion of $\left(x+\dfrac1x\right)^4$.
\begin{itemize}
\item[A)] $4$
\item[B)] $6$
\item[C)] $8$
\item[D)] $1$
\end{itemize}

\item In the expansion of $(1+x)^n$, the coefficient of $x^2$ is $21$ times the coefficient of $x^1$. Find $n$.
\begin{itemize}
\item[A)] $21$
\item[B)] $42$
\item[C)] $43$
\item[D)] $44$
\end{itemize}

\item The coefficients of $x^2$ and $x^3$ in the expansion of $(1+x)^n$ are equal. Find $n$.
\begin{itemize}
\item[A)] $4$
\item[B)] $5$
\item[C)] $6$
\item[D)] Cannot be determined.
\end{itemize}

\item Find the coefficient of $x^2$ in the expansion of $(1+2x+3x^2)^3$.
\begin{itemize}
\item[A)] $15$
\item[B)] $18$
\item[C)] $21$
\item[D)] $24$
\end{itemize}

\item The row-sums of Pascal's triangle are $1,2,4,8,16,\ldots$ Is this sequence of row-sums arithmetic, geometric, both, or neither? State the relevant ratio or difference.

\end{enumerate}

% ══════════════════════════════════════════════════════════════════════════════
\section*{Section C \quad Substitution \& Structure \hfill{\normalfont\small($\leq$90\,s each)}}
% ══════════════════════════════════════════════════════════════════════════════

\noindent\textit{Structure over computation --- what the binomial series is really telling you.}

\begin{enumerate}

\item Which of the following statements about partial sums of binomial coefficients is correct?
\begin{itemize}
\item[A)] $\binom n0+\binom n1+\cdots+\binom n{k-1}$ always equals $2^n$, for any $k\leq n$.
\item[B)] There is no simple closed form for a general partial sum of binomial coefficients, unlike the full row sum.
\item[C)] Partial sums of binomial coefficients are always arithmetic in $k$.
\item[D)] Partial sums of binomial coefficients are always geometric in $k$.
\end{itemize}

\item Find the coefficient of $x$ in the expression $(1+x)^0+(1+x)^1+\cdots+(1+x)^{40}$. \textit{\small(after TMUA 2019 Paper 1 Q3)}
\begin{itemize}
\item[A)] $40$
\item[B)] $41$
\item[C)] $820$
\item[D)] $1640$
\item[E)] $1681$
\end{itemize}

\item Find, in terms of $n$, the coefficient of $x^2$ in the expression $(1+x)^0+(1+x)^1+\cdots+(1+x)^n$. \textit{\small(after TMUA 2019 Paper 1 Q3)}
\begin{itemize}
\item[A)] $\dbinom n3$
\item[B)] $\dbinom{n+1}3$
\item[C)] $\dbinom{n+1}2$
\item[D)] $2^n-n-1$
\end{itemize}

\item In the expansion of $(1+x+x^2)^4$, what is the coefficient of $x^3$?
\begin{itemize}
\item[A)] $12$
\item[B)] $16$
\item[C)] $20$
\item[D)] $24$
\end{itemize}

\item Which of the following is the expansion of $(1-x)^{-1}$ as an infinite series in $x$?
\begin{itemize}
\item[A)] $1-x+x^2-x^3+\cdots$
\item[B)] $1+x+x^2+x^3+\cdots$
\item[C)] $1-2x+3x^2-4x^3+\cdots$
\item[D)] It cannot be expanded as a series.
\end{itemize}

\item The generalized binomial expansion gives $(1-x)^{-2}=\sum_{k=0}^\infty(k+1)x^k$. Find the coefficient of $x^3$.
\begin{itemize}
\item[A)] $3$
\item[B)] $4$
\item[C)] $6$
\item[D)] $10$
\end{itemize}

\item Which of the following is FALSE, for a positive integer $n$?
\begin{itemize}
\item[A)] $\displaystyle\sum_{k=0}^n\binom nk=2^n$
\item[B)] $\displaystyle\sum_{k=0}^n(-1)^k\binom nk=0$
\item[C)] $\dbinom nk=\dbinom n{n-k}$ for all $0\leq k\leq n$
\item[D)] $\displaystyle\sum_{k=0}^n k\binom nk=2^n$
\end{itemize}

\item The infinite series $1+x+x^2+x^3+\cdots$ equals $(1-x)^{-1}$. Using $(1-x)^{-3}=\sum_{k=0}^\infty\binom{k+2}{2}x^k$, find the coefficient of $x^5$ in $(1-x)^{-3}$.
\begin{itemize}
\item[A)] $\dbinom52$
\item[B)] $\dbinom62$
\item[C)] $\dbinom72$
\item[D)] $\dbinom82$
\end{itemize}

\end{enumerate}

% ══════════════════════════════════════════════════════════════════════════════
\section*{Section D \quad Challenge \hfill{\normalfont\small(TMUA / SMC difficulty)}}
% ══════════════════════════════════════════════════════════════════════════════

\noindent\textit{Insight over computation. Each question has a short, elegant route --- find it.}

\begin{enumerate}

\item Find the coefficient of $x^2y^3$ in the expansion of $(1+x+y^2)^5$. \textit{\small(after TMUA 2020 Paper 1 Q13)}
\begin{itemize}
\item[A)] $0$
\item[B)] $10$
\item[C)] $20$
\item[D)] $30$
\item[E)] $60$
\end{itemize}

\item The coefficient of $x^2$ in $(1+x+2x^2)^3$ equals the coefficient of $x^2$ in $(1-ax)^3$. Find the possible value(s) of $a$. \textit{\small(after TMUA Practice Paper 1 Q19)}
\begin{itemize}
\item[A)] $a=\pm\sqrt3$
\item[B)] $a=\pm3$
\item[C)] $a=\sqrt3$ only
\item[D)] $a=3$ only
\end{itemize}

\item Prove that $\displaystyle\sum_{k=0}^n k\binom nk=n\cdot2^{n-1}$ for every positive integer $n$, by differentiating the binomial expansion of $(1+x)^n$.

\item Expanding $(1+x)^n(1+x)^m$ two different ways --- first by combining the factors as $(1+x)^{n+m}$, then by multiplying the two separate expansions together and comparing the coefficient of $x^k$ on both sides --- produces which identity?
\begin{itemize}
\item[A)] Vandermonde's identity: $\displaystyle\sum_{i=0}^k\binom ni\binom m{k-i}=\binom{n+m}k$.
\item[B)] $\dbinom nk+\dbinom mk=\dbinom{n+m}k$.
\item[C)] $\dbinom nk\dbinom mk=\dbinom{n+m}k$.
\item[D)] No identity results --- the two sides are unrelated.
\end{itemize}

\item The coefficients of $(1-x)^{-1}$ are all $1$ (constant in $k$). The coefficients of $(1-x)^{-2}$, namely $(k+1)$, form an arithmetic sequence in $k$. What kind of sequence, as a function of $k$, do the coefficients of $(1-x)^{-3}=\sum_{k=0}^\infty\binom{k+2}2x^k$ form?
\begin{itemize}
\item[A)] Arithmetic.
\item[B)] Geometric.
\item[C)] Quadratic in $k$ (neither arithmetic nor geometric).
\item[D)] Constant.
\end{itemize}

\end{enumerate}

\SpeedClosing{``Pascal's triangle has been doing your multiplication for four hundred years.
Let it.''}

\end{document}
